\chapter*{CAPÍTULO VIII: CONCLUSIONES, RECOMENDACIONES Y BIBLIOGRAFÍA}
\addcontentsline{toc}{chapter}{CAPÍTULO VIII: CONCLUSIONES, RECOMENDACIONES Y BIBLIOGRAFÍA}

\section{CONCLUSIONES TÉCNICAS}

\begin{enumerate}
  \item \textbf{Suficiencia y Viabilidad Energética:} La Minicentral Hidroeléctrica Limbani (MCH Limbani), equipada con una turbina tipo Michell-Banki de flujo transversal y generador síncrono de $200.0\text{ kW}$ ($250.0\text{ kVA}$) a $380\text{ V}$, es plenamente capaz de abastecer la totalidad de la demanda eléctrica del casco urbano consolidado de Limbani. La máxima demanda coincidente calculada para el horizonte de diseño es de $\mathbf{181.80\text{ kW}}$ ($\mathbf{202.00\text{ kVA}}$ a $\cos\phi = 0.90$), lo cual representa un factor de utilización del $\mathbf{90.9\%}$ y deja un margen de reserva operativa del $\mathbf{9.1\%}$ ($18.2\text{ kW}$) para contingencias y pequeños talleres locales.
  \item \textbf{Calidad de Suministro y Perfiles de Tensión:} El diseño electromecánico adoptado garantiza una excelente calidad de tensión en todos los nodos de consumo, superando con holgura los estándares de la Norma Técnica de Calidad de los Servicios Eléctricos (NTCSE):
    \begin{itemize}
      \item En la línea primaria de 10.0 kV ($2.15\text{ km}$, conductor AAAC $35\text{ mm}^2$), la caída de tensión en hora de punta es de apenas $\mathbf{0.504\%}$ ($\Delta V = 50.41\text{ V}$), muy por debajo del límite regulatorio del $5.0\%$.
      \item En la red de baja tensión (380/220 V con cables autoportantes CAAI), la caída de tensión acumulada máxima en el punto más alejado de la SED-02 es de $\mathbf{3.82\%}$, cumpliendo con el valor admisible de $\le 5.0\%$.
    \end{itemize}
  \item \textbf{Alta Eficiencia y Bajas Pérdidas Técnicas:} Las pérdidas por efecto Joule en los $2.15\text{ km}$ de línea primaria en 10 kV ascienden a únicamente $\mathbf{1.05\text{ kW}}$, lo cual representa un insignificante $\mathbf{0.52\%}$ de la energía generada. Esto demuestra la conveniencia técnica de elevar la tensión a 10 kV mediante el transformador elevador de $250\text{ kVA}$ ($0.38/10\text{ kV}$) ubicado en la casa de máquinas, en lugar de realizar una distribución directa en baja tensión que habría resultado inviable por caída de tensión y costo de conductores.
  \item \textbf{Adaptación Climática y Rigor Normativo:} Todo el dimensionamiento de aislamiento, distancias mínimas de seguridad (DMS) y esfuerzos mecánicos en estructuras de concreto armado centrifugado (CAC) ha sido corregido aplicando el factor de derrateo por altitud $k_a = 1.28$ (correspondiente a $3,050\text{ m.s.n.m.}$) y las hipótesis de viento máximo de $90\text{ km/h}$ y temperatura extrema de $-5^\circ\text{C}$ de la Zona 2 (Sierra y Altiplano) del Código Nacional de Electricidad - Suministro 2011.
  \item \textbf{Seguridad de las Personas e Instalaciones:} El sistema de puesta a tierra calculado mediante el modelo de Dwight con dos jabalinas de cobre y tratamiento de bentonita sódica garantiza una resistencia de difusión de $\mathbf{8.14\ \Omega}$ en las tres subestaciones biposte, cumpliendo holgadamente el valor exigido por norma ($\le 10.0\ \Omega$).
  \item \textbf{Presupuesto y Plazo de Ejecución:} El presupuesto estimado referencial de obra asciende a $\mathbf{S/.\ 1,305,876.26}$ (incluyendo gastos generales, utilidad e IGV), con un plazo de ejecución física de $\mathbf{120\text{ días calendario}}$ (4 meses).
\end{enumerate}

\section{RECOMENDACIONES DE OPERACIÓN Y MANTENIMIENTO}

\begin{enumerate}
  \item \textbf{Mantenimiento Preventivo de la Turbina Michell-Banki:} Se recomienda establecer un programa riguroso de inspección trimestral del rodete y los álabes directrices de la turbina cross-flow, especialmente después de las temporadas de lluvias intensas (enero a marzo), a fin de evaluar y corregir el desgaste erosivo producido por sedimentos silíceos arrastrados por las aguas del Río Limbani.
  \item \textbf{Monitoreo Hidrológico en Época de Estiaje:} Durante los meses de mínimo estiaje (julio y agosto), se debe monitorear el caudal entrante al desarenador y cámara de carga de la MCH. En caso de que el caudal baje por debajo de los requerimientos de los 200 kW, se recomienda implementar un plan de gestión de demanda diurna y nocturna, priorizando los servicios esenciales (salud, alumbrado y bombeo de agua).
  \item \textbf{Inspección y Medición de Puestas a Tierra:} Semestralmente y antes del inicio de la temporada de tormentas eléctricas (octubre), se deberán medir las resistencias de los pozos a tierra de las tres subestaciones biposte y de los terminales de baja tensión con telurómetro digital, aplicando reactivación química con bentonita y sales electrolíticas si la resistencia superase los $10\ \Omega$.
  \item \textbf{Desbroce de Servidumbres de Línea en MT:} Mantener despejada la faja de servidumbre de $6.0\text{ m}$ (3.0 m a cada lado del eje de la línea de 10 kV), realizando podas de árboles de eucalipto y arbustos que puedan aproximarse a menos de $2.5\text{ m}$ de los conductores bajo ráfagas de viento.
  \item \textbf{Termografía en Transformadores y Tableros:} Realizar inspecciones termográficas anuales en los bornes de media tensión, bornes secundarios y barras de distribución de los tableros de las subestaciones SED-01, SED-02 y SED-03 para detectar a tiempo falsos contactos o desbalances térmicos.
\end{enumerate}

\section{REFERENCIAS BIBLIOGRÁFICAS Y BASE NORMATIVA}

\begin{enumerate}[label={[\arabic*]}]
  \item \textbf{Ministerio de Energía y Minas del Perú (MEM):} \textit{Código Nacional de Electricidad - Suministro 2011}, aprobado mediante Resolución Ministerial N.° 366-2011-MEM/DM.
  \item \textbf{Dirección General de Electricidad (DGE/MEM):} \textit{Bases para el Diseño de Líneas Primarias en Redes de Distribución}, aprobada por Resolución Directoral N.° 018-2003-EM/DGE.
  \item \textbf{Dirección General de Electricidad (DGE/MEM):} \textit{Bases para el Diseño de Redes Secundarias con Conductores Autoportantes}, aprobada por Resolución Directoral N.° 016-2003-EM/DGE.
  \item \textbf{Dirección General de Electricidad (DGE/MEM):} \textit{Especificaciones Técnicas para el Montaje de Líneas Aéreas y Redes de Distribución}, aprobada por Resolución Directoral N.° 026-2003-EM/DGE.
  \item \textbf{Ministerio de Energía y Minas del Perú:} \textit{Norma Técnica de Calidad de los Servicios Eléctricos (NTCSE)}, aprobada por Decreto Supremo N.° 020-97-EM y sus modificatorias.
  \item \textbf{Ministerio de Energía y Minas del Perú:} \textit{Reglamento de Seguridad y Salud en el Trabajo con Electricidad (RESESATE)}, aprobado por Resolución Ministerial N.° 111-2013-MEM/DM.
  \item \textbf{Instituto Nacional de Calidad (INACAL):} \textit{NTP 370.258: Conductores de aleación de aluminio para líneas aéreas (AAAC)}.
  \item \textbf{Instituto Nacional de Calidad (INACAL):} \textit{NTP 370.254: Cables concéntricos trenzados para distribución aérea (CAAI)}.
  \item \textbf{Instituto Nacional de Calidad (INACAL):} \textit{NTP 339.027: Postes de hormigón (concreto) armado para líneas aéreas}.
  \item \textbf{Instituto Nacional de Calidad (INACAL):} \textit{NTP 370.002: Transformadores de potencia. Especificaciones y pruebas}.
  \item \textbf{International Electrotechnical Commission (IEC):} \textit{IEC 60076: Power Transformers}; \textit{IEC 60502: Power cables with extruded insulation}; \textit{IEC 61089: Round wire concentric lay overhead electrical stranded conductors}.
  \item \textbf{Municipalidad Distrital de Limbani:} \textit{Catastro Urbano Oficial y Base Georreferenciada WGS 84 Zona 19 Sur}, 2026.
\end{enumerate}
